\documentclass[11pt]{article}

\usepackage[margin=1in]{geometry}
\usepackage{amsmath,amssymb}
\usepackage{booktabs}
\usepackage{hyperref}
\usepackage{microtype}
\setlength{\emergencystretch}{2em}

\newcommand{\UM}{U_{\mathrm{M}}}
\newcommand{\Usc}{U_{\mathrm{sc}}}
\newcommand{\Urel}{U_{\mathrm{rel}}}
\newcommand{\Ebeam}{E_{\mathrm{beam}}}
\newcommand{\Eemit}{E_{\mathrm{emit}}}
\newcommand{\code}[1]{\texttt{#1}}

\title{Prospective Restricted-Domain Validation of an Upward-Beam\\
Spacecraft-Relative Potential Proxy in Lunar Prospector ER Data}
\author{ESPG Lunar Potential Map Project}
\date{11 July 2026}

\begin{document}
\maketitle

\begin{abstract}
An upward field-aligned electron feature provides the most stable voltage-scale
observable recovered in our clean-room analysis of the Lunar Prospector
Electron Reflectometer benchmark.  A subsequent full-domain simulation study
found calibrated held-out centroid inference at 150 and 350~eV but failed at
65 and 800~eV and exceeded its beam-absent false-detection gate.  Because the
candidate 100--500~eV operating window was proposed after those outcomes were
known, it requires a new prospective experiment.  We preregister a
polarity-aware, sector-level validation in which continuum structure is
independent of injected beam energy, calibration and confirmation use disjoint
instrument supports, two normalizations are derived from a common latent
observation, and response, morphology, temporal, semi-synthetic, null, and
out-of-window challenges are explicit.  The primary matrix contains 4800
moderate-to-strong beams over twelve energy-amplitude cells, accompanied by 600
weak beams, 1200 semi-synthetic injections, 3000 nulls, and 1000 out-of-window
classification controls.  Retention is co-primary with conditional interval
coverage, and every cell must pass.  A full pass would validate selective
measurement of a pre-response beam centroid over the tested envelope and open
a hand-audited one-day real-data study.  It would not validate absolute lunar
surface potential, spacecraft potential, loss-cone corroboration, or a global
map.  No simulation result had been generated when this design was written.
\end{abstract}

\section{Scientific Background}

For a surface-emitted electron transported electrostatically to the spacecraft,
energy conservation gives, in electron-volt/volt numerical units,
\begin{equation}
  \Ebeam \simeq \Eemit-\Urel+\delta_{\mathrm{transport}},
  \qquad \Urel=\UM-\Usc.
  \label{eq:beam}
\end{equation}
The commonly used proxy $\Urel\simeq-\Ebeam$ sets the initial secondary energy
and transport discrepancy to zero.  This is often a useful leading-order
scale, but it is not an identity.  Absolute lunar potential additionally
requires an independent $\Usc$ estimate:
\begin{equation}
  \UM=\Urel+\Usc.
\end{equation}

Experiment~004 found a robust response-aware beam scale near 145--160~eV in the
historical benchmark but did not find stable independent voltage information
in the loss-cone boundary.  Experiment~005 therefore tested beam-centroid
inference directly.  Its smooth-continuum estimator achieved calibrated
held-out results in both 150 and 350~eV cells, but weak low-energy and both
high-energy conditions failed; beam-absent acceptance was 6\% against a frozen
5\% point-estimate gate.  The complete 20--2000~eV hypothesis failed.

The observed middle-energy success motivates a restricted hypothesis.  It does
not confirm one.  The present experiment is designed so that the new window can
be rejected without another retrospective rescue.

\section{Lessons From The Parent Experiment}

Several limitations became visible only when the complete Experiment~005 record
was reconsidered.

First, the synthetic loss-cone boundary was generated using
$\Urel=-\Ebeam$.  The nuisance continuum therefore contained indirect
information about the injected answer even though the fitted continuum itself
had no voltage equation.  The new primary generator samples continuum
parameters independently of beam energy.

Second, both normalization branches received the same already-normalized
synthetic observation.  Their agreement was automatic.  The new experiment
starts from one latent native-sector flux realization and derives incident-mean
and conjugate-sector normalizations independently.

Third, pitch was computed without establishing which magnetic-field direction
connected to the surface.  Here surface intersection and polarity are required
before defining 180 degrees as surface-leaving.  Disconnected geometry is
excluded rather than assigned a directional interpretation.

Fourth, four support spectra were shared by calibration and validation,
coverage was conditional on selection without a strong retention requirement,
and the null point estimate was underpowered as an upper-bound statement.  The
new design uses disjoint outcome-blind support sets, makes retention co-primary,
and uses exact binomial confidence limits.

\section{Statistical Estimand}

Let $S=1$ denote passage of every frozen geometry, numerical, normalization,
stability, and interval-resolution rule.  The primary estimands are
\begin{align}
  r_j &= \Pr(S=1\mid j),\\
  b_j &= \operatorname{median}
  \left[\frac{\widehat E_b-E_b}{0.48E_b}\middle|S=1,j\right],\\
  c_{q,j} &= \Pr(E_b\in I_q\mid S=1,j),\qquad q\in\{0.68,0.95\},
\end{align}
for each prespecified energy-amplitude cell $j$.  Reporting coverage only
conditional on $S=1$ is legitimate only when $r_j$ is reported and gated.
Every simulation, including a crash or unresolved fit, remains in the
retention denominator.

For a beam-absent family $k$, define
\begin{equation}
  p_{0,k}=\Pr\left(S=1,\ 100\leq\widehat E_b\leq500
  \mid\text{no beam},k\right).
\end{equation}
The gate is not $\widehat p_{0,k}\leq0.05$.  It requires the
Bonferroni-adjusted one-sided familywise-95\% Clopper--Pearson upper bound to be
at most 0.05.  With three families of 1000 controls, no family may contain more
than 35 false in-window acceptances.

\section{Estimator And Reporting Domain}

The observed support and optimizer remain broad, 20--2000~eV.  The reporting
window, 100--500~eV, is applied after fitting.  This separation prevents a null
or out-of-window feature from being forced against a newly narrowed optimizer
bound.  Confidence intervals are not clipped at 100 or 500~eV, and base QC is
computed without conditioning on whether the fitted point falls inside the
window.

The model is
\begin{equation}
  R(E,\alpha)=C(E,\alpha)+A
  G_E(E;E_b)G_\alpha(\alpha;180^\circ,\sigma_\alpha).
  \label{eq:model}
\end{equation}
The pitch-logistic continuum $C$ has cone level, contrast, boundary, and width
varying quadratically in normalized log energy.  It includes no adiabatic
boundary equation.  The nominal energy response has
$\mathrm{FWHM}/E=0.48$ and is integrated over channel support.  The fitter uses
the frozen unweighted log-residual objective, deterministic multistart, and
quality checks inherited in concept from Experiment~005.

\section{Support And Information Barriers}

Twenty-four supports are selected from the 398-sweep parent interval using only
time, structural validity, connection, pitch coverage, native energy alignment,
valid-sector counts, and magnetic variation.  Known Experiment~005 supports
and the Halekas neighborhood cannot enter confirmation.  Twelve supports are
assigned to calibration and twelve disjoint supports to confirmation.

The experiment has four barriers.  First, support selection and geometry are
committed.  Second, implementation verification and fresh calibration are run.
Third, code, corrections, seeds, checksums, and decision logic are frozen.
Fourth, confirmation runs once without incremental outcome inspection.  Any
method change after that point requires unused draws and a new experiment.

\section{Calibration And Confirmation}

Calibration uses 3000 beams at 100, 140, 220, 340, and 500~eV.  One
deterministic half estimates median bias and the other estimates rank-based 68\%
and 95\% interval multipliers.  No Experiment~005 row or correction value is
reused.

Confirmation crosses six energy strata---the exact endpoints and four
log-uniform interior bands---with amplitudes 4--8 and 8--16.  Four hundred
cases in each cell give 4800 primary beams.  Pitch width, noise, support,
continuum, response, and channel phase are balanced independently within every
cell.  Six hundred amplitude-1--4 beams characterize, but cannot expand, the
primary claim.

The primary generator includes independent smooth, double-transition,
asymmetric, and correlated energy-pitch backgrounds.  Challenges include
heteroskedastic and correlated noise, channel integration, response-width and
energy-scale stress, log-energy and asymmetric response, beam-axis offsets,
double peaks, missing channels, and temporal changes during the sequential
energy sweep.  A further 1200 injections use held-out real residual templates.

Three thousand beam-absent controls cover matched, misspecified, and adversarial
families.  The adversarial family includes energy-only and pitch-only
enhancements, single-row spikes, incident-side features, and temporal ramps or
pulses.  One thousand true out-of-window beams test leakage into the reporting
window; the 500 cases on each side may produce no more than 15 final-valid
in-window reports under their adjusted exact-binomial gate.

\section{Decision Rule}

Every one of twelve primary cells must retain at least 200 of 400 attempts.
Among retained cases, the bootstrap upper confidence bound for median absolute
error must not exceed 0.10 response FWHM, the corresponding bound for the 90th
percentile must not exceed 0.25 FWHM, and the signed-median interval must lie
within $\pm0.05$ FWHM.  Empirical 68\% and 95\% coverage must lie in
0.60--0.76 and 0.90--0.99, respectively, with each nominal level contained in
its exact 95\% binomial interval.  Every retained 95\% interval must be no wider
than one response FWHM.  Pooled background-family gates, semi-synthetic gates,
null gates, and out-of-window leakage gates must also pass.

These requirements form an intersection claim: the complete window passes only
when all required components pass.  Failure at an edge rejects the full window.
Failure of null control may leave centroid recovery interesting for an
independently established beam, but it does not validate automated detection.
Failure of the semi-synthetic layer limits any success to the simulation system.

\section{Expected Outcomes}

Three outcomes are scientifically plausible.  First, the middle-energy success
may generalize across new supports and energy locations, including the proposed
edges.  This would justify a one-day, hand-audited observable study.  Second,
the central regime may remain valid while one or both edges fail.  This rejects
the 100--500~eV claim and supplies an exploratory hypothesis for a later
experiment; it is not a partial pass.  Third, response, polarity, normalization,
temporal, or semi-synthetic challenges may reveal that the earlier success was
too simulator-specific.  That negative result would prevent a misleading real
survey and identify the missing physics or observation model.

\section{Interpretation Limits}

A full pass validates a selective measurement procedure for $\Ebeam$ over its
tested amplitude, morphology, support, and discrepancy envelope.  It does not
show that every accepted real feature is a lunar secondary, explain beam
absence, measure real-data bias without truth, validate $\Usc$, recover $\UM$,
provide independent loss-cone corroboration, establish mission-wide
representativeness, or reverse the global reconstruction's identifiability
limit.

The real-data stage, if opened, is therefore a residual and prevalence audit.
Its purpose is to learn whether the validated observable occurs in actual LP ER
sweeps under the same quality rules, not to jump directly from a successful
simulation to a lunar potential map.

\section{Implementation Certification Result}

Before calibration, we froze a 512-case implementation corpus containing 96
matched analytic cases, 160 generic discrepancy cases, 64 null or adversarial
cases, and 192 cases on the twelve calibration supports.  Each case used the
same twelve truth-blind starts.  Complete SciPy multistart fitting defined the
reference; float64 projected Levenberg--Marquardt GPU candidates were ranked by
the canonical NumPy objective and polished on the CPU.  Certification required
100\% reference-basin recall for some $K\leq6$ and final objective, modeled
signal, and identifiable-center parity.

The GPU-candidate procedure failed this gate.  Rank-1 recall was 82.3\% in the
matched analytic arm, 92.5\% under generic discrepancy, 88.0\% on calibration
supports, and 62.5\% in the null/adversarial arm.  Overall recall increased
only from 85.2\% at $K=1$ to 86.1\% at $K=12$.  Thus the failure was not chiefly
an ordering problem among multiple useful GPU basins: in many cases the
candidate set contained no eligible trajectory in the CPU-reference basin.
Only 337 of 512 cases supplied at least six converged GPU candidates, and the
fail-closed hybrid invoked complete CPU fitting in 226 cases.

Seven deliberately non-identifiable controls---one flat-pitch, one
single-energy, and five zero-dynamic-range cases---had no successful finite CPU
reference trajectory.  They remain failures in the denominator.  The result
does not challenge the already verified float64 CPU/CUDA forward equivalence,
and it does not test the beam-proxy hypothesis.  It rejects only replacement or
pruning of complete CPU multistart by the current projected-LM candidates.

Consequently, fresh calibration will use the complete deterministic twelve-start
SciPy fitter.  Retaining this slower path prevents an implementation shortcut
from creating energy- or support-dependent selection before scientific
calibration.  Redesigning the GPU optimizer would require a new, independently
frozen certification corpus and is not part of the present experiment.

\section{Current Status}

Support selection and implementation certification are complete.  No
confirmation flux has been opened.  Fresh calibration, correction freeze, and
one-shot confirmation remain to be run under the preregistered barriers.

\end{document}
